\begin{table}[htbp]
\centering
\small
\caption{Состав корпуса и условия производства текста}
\label{tab:1}
\begin{tabular}{llll}
\toprule
Часть корпуса & Условие производства & Категория & Документов \\
\midrule
машинная & канал генерации & real\_claude & 270 \\
машинная & канал генерации & deepseek\_pro & 270 \\
машинная & канал генерации & gpt & 270 \\
машинная & канал генерации & nemotron & 269 \\
машинная & режим задания & P1 & 360 \\
машинная & режим задания & P2 & 360 \\
машинная & режим задания & P3 & 359 \\
машинная & повтор & 1 & 540 \\
машинная & повтор & 2 & 539 \\
человеческая & уровень 2 & Шаблонные материалы, seo & 231 \\
человеческая & уровень 0 & Свободные авторские тексты, prose & 180 \\
человеческая & уровень 3 & Стандартизованная проза, science & 167 \\
человеческая & уровень 1 & Публицистика, news & 152 \\
человеческая & уровень не присваивается & Тексты языкового контакта, translation & 73 \\
обе & итог & аналитический корпус & 1882 \\
\bottomrule
\end{tabular}
\begin{minipage}{\linewidth}\footnotesize\vspace{2pt}
\textit{Примечание.} Единица анализа — документ. Знаменатель: 1882 документов аналитического корпуса, 1079 машинных и 803 человеческих. Данные: реестр корпуса, профиль препроцессинга prep-v5. Приведены счётчики документов; оценок и интервалов в таблице нет. Шкала: шкалы нет, таблица описывает состав. Сокращения: P1, P2 и P3 — режимы задания: P1 — нейтральное техническое задание, P2 — жёсткое редакционное ТЗ, P3 — требование избегать типичных признаков машинного текста; уровень — степень регламентированности источника, присвоенная по источнику до чтения текста. Пропуски: страта translation уровень не получает: у переводов регламент задаёт исходный текст, а не редакция. Поправка на множественные сравнения не применялась.
\end{minipage}
\end{table}
